\begin{table}[h]
\centering
\footnotesize
\setlength{\tabcolsep}{3pt}
\begin{adjustbox}{max width=\linewidth}
\begin{tabular}{lcc}
\toprule
& Chain 1 against chain 3 & Chain 1 against chain 2 \\
\midrule
Items & 310 & 310 \\
Difference of $D$ [CI] & .001 [$-.005$, .007] & .000 [$-.004$, .004] \\
Pearson [CI] & .614 [.412, .770] & .404 [.292, .510] \\
Spearman [CI] & .667 [.578, .744] & .368 [.258, .473] \\
Intraclass correlation [CI] & .614 [.408, .769] & .404 [.291, .509] \\
Share of variance between items & .608 & .404 \\
\bottomrule
\end{tabular}
\end{adjustbox}
\caption{Reliability of the item differences behind $D$ for Haiku 5.5 off. Chain 1 sends the fields on half of the items, chain 2 repeats the request for a first answer and sends the fields on every item, and chain 3 sends no fields. Chain 1 is compared with chain 3 on the half where chain 1 sent no fields and with chain 2 on the half where chain 1 sent the fields, so that the two chains compared in a column sent the same request. $D$ is 0.084 in chain 1, 0.040 in chain 2 and 0.128 in chain 3 over all items, and chain 3 minus chain 2 is 0.089 (95 percent interval 0.084 to 0.094). The mixed model has random intercepts for events and items and converged with a warning from the optimiser. The intervals of the Spearman correlations follow \citet{bonett2000}.}
\label{tab:reliability}
\end{table}
